\documentclass{report}
\usepackage{amsmath,amssymb,graphicx}
\setlength{\parindent}{0mm}
\setlength{\parskip}{1em}
\begin{document}
\begin{center}

\rule{6in}{1pt} \
{\large
Ilya Lashuk
\\ {\bf
Preconditioned eigensolvers in Hypre and PETSc
}}

Department of Mathematics \\
University of Colorado \\
P.O. Box 173364, Campus Box 170 \\
Denver CO 80217-3364 
\\ {\tt
ilashuk@math.cudenver.edu
}
\\
Merico Argentati,
Evgueni Ovtchinnikov,
Andrew Knyazev

\includegraphics[width=60mm]{lashuk}
\end{center}

We present preliminary results of an ongoing project to develop codes
of the Locally Optimal Block Preconditioned Conjugate Gradient (LOBPCG)
method for symmetric eigenvalue problems for Hypre and PETSc software
packages. Hypre and PETSc packages provide high quality multigrid and
domain decomposition preconditioning on parallel clusters with
distributed or shared memory architecture. The LOBPCG method, suggested
and developed by Andrew Knyazev [1] in the past decade, recently
attracts an increasing attention as a potential alternative to the
shift-and-invert Lanczos and preconditioned Davidson methods due to its
simplicity robustness and fast convergence. Several MATLAB, C, C++ and
FORTRAN implementations of the LOBPCG are developed by different
groups, e. g., for such applications areas as structured mechanics and
electronic structure calculations. However, the only publicly available
at present LOBPCG implementation remains our implementation for Hypre,
which is already included in Hypre 1.8.2 release. We describe the
current state of the LOBPCG software for Hypre and PETSc, developed by
our group and demonstrate initial scalability results.

[1] A.V. Knyazev, {\em Toward the Optimal Preconditioned Eigensolver:
Locally Optimal Block Preconditioned Conjugate Gradient Method}.
SIAM Journal on Scientific Computing {\bf 23} (2001), 2, pp.517--541.

\end{document}
